\chapter{Introducción}

La computación concurrente y paralela constituye una de las principales estrategias para incrementar el desempeño de aplicaciones científicas y de análisis de datos, especialmente cuando el volumen de información y el costo computacional de los algoritmos superan las capacidades de ejecución secuencial tradicionales.

El presente documento describe detalladamente el diseño e implementación del proyecto \texttt{PDC\_FINAL}, desarrollado en el contexto académico de la asignatura Programación Concurrente y Distribuida.

El proyecto implementa un algoritmo de búsqueda aleatoria para maximizar la métrica Area Under the Curve (AUC), empleando múltiples paradigmas de paralelización con el propósito de analizar su comportamiento, rendimiento y escalabilidad.

Las implementaciones desarrolladas corresponden a:

\begin{itemize}
    \item Implementación secuencial en Python.
    \item Implementación multicore mediante \texttt{multiprocessing}.
    \item Implementación en memoria compartida usando OpenMP.
    \item Implementación distribuida utilizando MPI.
    \item Implementación acelerada por GPU mediante PyCUDA.
\end{itemize}